\documentclass{article}
\usepackage{graphicx} % Required for inserting images
\usepackage{amsmath}
\usepackage{float}
\usepackage{amssymb}
\title{Lesson 7 Homework - Proof Writing}
\author{William Wang}
\date{October 2025}

\begin{document}

\maketitle

\section{}
Since we know that $y = \arctan x \implies \tan y = x = \frac{x}{1}$, we know that we can construct a triangle with opposite side length $x$ and adjacent side length 1 and by the Pythagoras Theorem, we have the 
\[\text{Hypotenuse} = \sqrt{x^2+1}\]
\begin{figure}[H]
    \centering
    \includegraphics[width=0.75\linewidth]{q1.png}
    \caption{$y$ is the angle labelled}
\end{figure}
Then we have 
\[\sin y = \frac{\text{opposite}}{\text{hypotenuse}} = \frac{x}{\sqrt{x^2+1}}\]
By rationalizing the denominator, we have 
\[\frac{x\sqrt{x^2+1}}{x^2+1}\]
\section{}
Since we have 
\[3\cot(4k-\pi) = \sqrt{3} \implies \cot(4k-\pi) = \frac{\sqrt{3}}{3}\]
Since we have $\cot x = \frac{1}{\tan x}$, we have 
\[\tan(4k-\pi) = \frac{3}{\sqrt{3}}\]
Then let $\theta = 4k-\pi$, we can construct a triangle: 
\begin{figure}[H]
    \centering
    \includegraphics[width=0.75\linewidth]{q2.png}
    \caption{The constructed triangle (not to scale)}
\end{figure}
Notice we can multiply each side by $\frac{\sqrt{3}}{3}$ so that we side lengths $(1, \sqrt{3}, 2)$, notice this is a 30-60-90 triangle. Hence we know that the angle
\[\theta = \frac{\pi}{3}\implies 4k-\pi = \frac{\pi}{3} \implies k = \frac{\pi}{3}\]
However since $\cot$ is a periodic function and has a cycle of $\pi$, we know that this transformed function has cycle of 
\[\frac{\pi}{|B|} = \frac{\pi}{4}\]
Hence the set of solution is 
\[\{ \frac{\pi}{3} + \frac{\pi x} {4} \mid x \in \mathbb{Z}\}\]

\end{document}
